\stfhIndex{I129}{Zhongdenghuofeierfanshu}{中等霍费尔范数（medium Hofer norm）}
\section{霍费尔测地线与周期轨道}
\label{sec:hofer-geodesics-periodic-orbits}

本节研究哈密顿路径在霍费尔距离下的长度极小性。我们先沿用\personBialy{}与\personPolterovich{}
(BP94) 以及拉隆德与麦克达夫 (LM95b) 的做法，引入一个必要判据，再介绍厄舍 (Ush10a)
对所谓极小性猜想的证明。同时，我们也将依照厄舍的论证证明
$\operatorname{Ham}(M,\omega)$ 上 $\gamma$-范数的非退化性。

\subsection{准自治哈密顿量}

$H$ 的 $L^{(1,\infty)}$-范数
$\lVert H\rVert=\int_0^1(\max H_t-\min H_t)\,dt$ 可等同于路径
$\phi_H:t\mapsto\phi_H^t$ 的芬斯勒长度
\[
 \operatorname{leng}(\phi_H)=\int_0^1
 \bigl(\max_xH(t,\phi_H^t(x))-\min_xH(t,\phi_H^t(x))\bigr)\,dt,
\]
其中
\[
 T_{\mathrm{id}}\operatorname{Ham}(M,\omega)
 \cong C^\infty(M)/\mathbb R
\]
上的相伴巴拿赫范数正是
\[
 \lVert h\rVert=\operatorname{osc}(h)=\max h-\min h.
\]
映射 $t\mapsto H(t,\phi_H^t(\mathord\cdot))$ 可视为
$\operatorname{Ham}(M,\omega)$ 上路径 $t\mapsto\phi_H^t$ 的切向量。

\begin{definition}
\label{def:22-hofer-path-metric}
考虑定义在
$\mathcal P^{\mathrm{ham}}(\operatorname{Symp}(M,\omega),\mathrm{id})$ 上的度量
$d:\mathcal P^{\mathrm{ham}}(\operatorname{Symp}(M,\omega),\mathrm{id})
\to\mathbb R_+$：
\[
 d(\lambda,\mu):=\operatorname{leng}(\lambda^{-1}\circ\mu),
\]
其中 $\lambda^{-1}\circ\mu$ 是哈密顿路径
$t\in[0,1]\mapsto\lambda(t)^{-1}\mu(t)$。路径空间
$\mathcal P^{\mathrm{ham}}(\operatorname{Symp}(M,\omega),\mathrm{id})$
上的霍费尔拓扑即由此度量诱导。
\end{definition}

\begin{exercise}
\label{ex:22-hofer-quotient-topology}
考虑评价映射
\[
 \operatorname{ev}_1:
 \mathcal P^{\mathrm{ham}}(\operatorname{Symp}(M,\omega),\mathrm{id})
 \longrightarrow\operatorname{Ham}(M,\omega).
\]
假设 $M$ 紧。证明：上述 $\operatorname{Ham}(M,\omega)$ 上的霍费尔拓扑，是使评价
映射 $\lambda\mapsto\lambda(1)$ 连续的最强拓扑。
\end{exercise}

该拓扑的非平凡性来自霍费尔范数函数
$\lVert\mathord\cdot\rVert:\operatorname{Ham}(M,\omega)\to\mathbb R_+$
的非退化性。霍费尔 (H93) 证明：只要相应哈密顿方程不存在\emph{周期小于或等于一}的
非常值周期轨道，$\mathbb C^n$ 上任意紧支\emph{自治}哈密顿量的路径便是长度极小的。
(En00), (MSl01) 与 (Oh02)--(Oh05b) 在附加假设每个不动点处的线性化流不是过扭的
条件下推广了这个结果；亦即，该线性化流没有周期小于一的闭轨道。

\personBialy{}--\personPolterovich{}在 (BP94)、\personLalonde{}--\personMcDuff{}在 (LM95b) 中证明：任意长度极小
（分别地，局部长度极小）的哈密顿路径，都由\emph{准自治}（分别地，\emph{局部准自治}）哈密顿路径
生成。

\begin{definition}
\label{def:22-quasi-autonomous}
若存在两个在时间 $t\in[0,1]$ 内保持不动的点 $x_-,x_+\in M$，使得对所有
$t\in[0,1]$ 都有
\[
 H(t,x_-)=\min_xH(t,x),\qquad
 H(t,x_+)=\max_xH(t,x),
\]
\stfhSubIndex{I106}{Hamidunliang}{哈密顿量（Hamiltonian）}{Zhunzizhi}{准自治（quasi-autonomous）}
则称哈密顿量 $H$ 为\emph{准自治的}。
\end{definition}

下述定义一直用于研究哈密顿路径的长度极小性质。

\begin{definition}
\label{def:22-twisted-periodic-point}
设 $H:[0,1]\times M\to\mathbb R$ 是不必时间周期的哈密顿量，$\phi_H^t$ 是其哈密顿流。
\begin{enumerate}
 \item 若 $\phi_H^T(p)=p$，则称点 $p\in M$ 是\emph{时间-$T$ 周期点}。若轨道
 $t\in[0,T]\mapsto\phi_H^t(p)$ 可缩，则称之为\emph{可缩时间-$T$ 周期轨道}。
 \item 若 $H$ 在 $t\in[0,T]$ 内有不变临界点 $p$，且其线性化流
 \[
  d\phi_H^t(p):T_pM\longrightarrow T_pM,\qquad t\in[0,T],
 \]
 有周期小于或等于 $T$ 的闭轨道，则称 $p$ 作为时间-$T$ 周期轨道是\emph{过扭的}；
 否则称为\emph{欠扭的}。若线性化流此外只有原点作为不动点，则称该不动点为
 \emph{通有欠扭的}。
\end{enumerate}
\end{definition}

注意，通有欠扭性是哈密顿量 $H$ 的 $C^2$-稳定性质。

为引出后面的讨论，我们回忆\personUstilovsky{}--\personLalonde{}--\personMcDuff{}关于测地线稳定性的必要
条件，但略去证明。

\begin{theorem}[\personUstilovsky{} (Ust96)，\personLalonde{}--\personMcDuff{} (LM95b)]
\label{thm:22-stable-geodesic-fixed-points}
设 $\phi$ 是通有的，即其所有不动点均孤立。则从恒等映射到 $\phi$ 的任意稳定测地线
$\phi^t$（$0\leq t\leq1$），必有至少两个欠扭不动点。
\end{theorem}

下面的猜想由波尔捷罗维奇提出（(Po01) 猜想 12.6.D），有时称为极小性猜想。
（亦见 (Po01), (LM95b), (MSl01)。）

\begin{conjecture}[极小性猜想]
\label{conj:22-minimality}
任意不存在周期小于或等于一的可缩周期轨道的自治哈密顿路径，都在其端点固定的路径同伦类
中是霍费尔长度极小的。
\end{conjecture}

恩托夫 (En00) 的下述定义简化了术语。

\begin{definition}
\label{def:22-slow-hamiltonian}
\stfhSubIndex{I107}{Hamidunliang}{哈密顿量（Hamiltonian）}{Man}{慢（slow）}
若哈密顿量 $H$ 没有周期小于或等于一的可缩周期轨道，则称 $H$ 为\emph{慢的}。
\end{definition}

在相伴自治哈密顿路径的不动点线性化流满足各种条件时，多位作者证明了极小性猜想。
厄舍 (Ush10a) 在不对线性化流附加任何限制的情况下完整证明了此猜想。他使用
(Oh05b) 中的长度极小判据以及哈密顿量的一条逼近定理。（相关结果亦见 (Sch06)。）
下面两个小节将介绍厄舍的证明。

\subsection{\texorpdfstring{通过 $\rho(H;1)$ 的长度极小判据}{通过 rho(H;1) 的长度极小判据}}

我们用 (Oh05b) 中的谱不变量 $\rho(H;1)$ 描述哈密顿路径长度极小性质的一个简单判据。
该判据类似于 (H93) 以及 $\mathbb C^n$ 情形的 (BP94) 中所用的判据。

为描述这一判据，把~\eqref{eq:22-hofer-lengths} 中定义的范数分解为
\[
 \lVert H\rVert=E^-(H)+E^+(H)
\]
会很有用。

\begin{proposition}
\label{prop:22-negative-length-minimizing}
设 $G:[0,1]\times M\to\mathbb R$ 是满足
\begin{equation}
 \rho(G;1)=E^-(G)
 \label{eq:22-negative-length-criterion}
\end{equation}
的任意哈密顿量。则 $G$ 在其端点固定的同伦类中使负霍费尔长度极小。特别地，若此外
$\rho(\overline G)=E^-(\overline G)=E^+(G)$，则 $G$ 必为准自治的。
\footnote{译者注：原书命题首项结论印作“Then $H$”，但命题只给定了 $G$；译文据证明
中的 $F\sim G$ 改作 $G$。}
\end{proposition}

\begin{proof}
设 $F$ 是满足 $F\sim G$ 的任意哈密顿量。由~\eqref{eq:22-negative-length-criterion}、
(21.6.39)（取 $a=1$）以及~\eqref{eq:22-classical-hofer-bounds}，依次得到
\[
 E^-(G)=\rho(G;1)=\rho(F;1)\leq E^-(F).
\]
这证明了第一项陈述。第二项陈述于是由\personBialy{}--\personPolterovich{}、\personUstilovsky{}以及
拉隆德--麦克达夫的极小性判据得出。
\end{proof}

另一方面，\emph{若 $G$ 是一周期的}，就可以考虑相伴作用量泛函 $\mathcal A_G$。对于准自治
周期哈密顿量 $G$，$\mathcal A_G$ 有两个显然的临界值：
\[
 \mathcal A_G([x_-,\widetilde x_-])
 =\int_0^1-G(t,x_-)\,dt,
 \qquad
 \mathcal A_G([x_+,\widetilde x_+])
 =\int_0^1-G(t,x_+)\,dt.
\]
要点在于，这些作用量值分别与 $\lVert H\rVert$ 的负、正部分
\[
 E^-(G)=\int_0^1-\min G_t\,dt,
 \qquad
 E^+(G)=\int_0^1\max G_t\,dt
\]
一致。


\subsection{基本弗洛尔闭链的构造}

现在引入同调本质性的一个基本概念，即\emph{在柯西--黎曼流下的不可下推性质}。
这是\emph{链层弗洛尔理论}的核心，特别是谱不变量的构造 (Oh05c) 及其应用
(Oh05b, Oh05d, Oh09a) 的核心。这个概念最早在 (Oh97b) 关于余切丛中正合拉格朗日
子流形的语境下隐含地引入弗洛尔理论；随后在一般哈密顿微分同胚的语境下，作为
紧弗洛尔闭链的概念明确出现于 (Oh02, Oh09a)。

由定义 21.5.2 中紧弗洛尔闭链的定义，若量子同调类满足 $[\alpha]=a^\flat$，则
$(H,J)$ 的任意紧弗洛尔闭链都实现数值 $\rho(H;a)$，即
$\lambda_{(H,J)}(\alpha)=\rho(H;a)$。证明这种闭链存在并非易事。下述定理在 (Oh09a)
中由几何连续性方法证明，在 (Ush08) 中由纯代数方法证明。

\begin{theorem}[吴龙根 (Oh09a)，厄舍 (Ush08)]
\label{thm:22-tight-floer-cycle-exists}
对任意 $0\neq a\in QH^*(M)$，都存在一个与 $a$ 相伴的紧弗洛尔闭链。
\end{theorem}

此定理的一般证明相当微妙（除非 $(M,\omega)$ 是有理的），超出本书范围。因此请读者
参阅原始论文。

与命题~\ref{prop:22-negative-length-minimizing} 中的长度极小判据有关，对满足慢性条件的
准自治哈密顿量 $G$，关键在于构造一个与 $a=1$ 相伴、层级
$\lambda_{(G,J)}(\alpha)=E^-(G)$ 的紧弗洛尔闭链。

然而，一般哈密顿量的弗洛尔闭链定义因其表述方式而相当抽象复杂，故通常不易用于实际
应用。另一方面，给定 $(M,\omega)$ 的基本闭链所对应的弗洛尔闭链，可以通过弗洛尔链
映射传递莫尔斯基本闭链来系统研究。

\begin{definition}
\label{def:22-fundamental-floer-cycle}
若弗洛尔闭链 $\alpha$ 的同调类满足 $[\alpha]=1^\flat$，则称 $\alpha$ 为\emph{基本闭链}。
\end{definition}

类似地，莫尔斯基本闭链在莫尔斯理论中定义。当哈密顿路径满足准自治性质之类的特定
性质时，我们希望在研究这些传递所得的闭链时利用这一几何性质。需要强调的是，一般而言，
\emph{这个过程只在弗洛尔链层面而非临界点层面良定义}。这是因为在弗洛尔链映射下，一个临界点
被传递为一条弗洛尔链，而不仅是另一个临界点。

选择莫尔斯函数 $f$，并取\emph{充分小的} $\epsilon>0$ \emph{使定理 19.8.1 成立}；再取
$-\epsilon f$ 的莫尔斯基本闭链 $\alpha$。于是也可把 $\alpha$ 视为
$\epsilon f$ 的弗洛尔闭链。通过弗洛尔链映射 $h_{\mathcal L}$ 传递 $\alpha$，定义
$H$ 的基本弗洛尔闭链
\[
 \alpha_H:=h_{\mathcal L}(\alpha)\in CF(H),
\]
其中 $h_{\mathcal L}$ 是任意连接 $\epsilon f$ 与 $H$ 的给定同伦 $\mathcal L$ 上的
弗洛尔链映射。关于极小性猜想，典范给出的线性路径
\[
 \mathcal L:s\longmapsto(1-s)\epsilon f+sH
\]
起着重要作用。还要注意，这条闭链依赖于莫尔斯函数 $f$ 的选择。一般而言，即使 $H$
是准自治的，也不能期待这条闭链是紧的。

\stfhIndex{I116}{Keerman-Lalongde}{克尔曼--拉隆德（Kerman--Lalonde）}
在这方面，克尔曼--拉隆德 (KL03) 证明：在下面两项要求下，
\begin{enumerate}
 \item 准自治哈密顿量 $H$ 有唯一的非退化整体极小点 $x_-$，且该点对所有
 $t\in[0,1]$ 均欠扭；
 \item 出现在上述线性路径 $\mathcal L$ 中的莫尔斯函数 $f$
 \begin{itemize}
  \item 以 $x_-$ 为整体极小点；
  \item 且对所有其他临界点 $x_j$ 满足
  \begin{equation}
   f(x_-)=0,\qquad f(x_-)<f(x_j),
   \label{eq:22-morse-minimum-condition}
  \end{equation}
 \end{itemize}
\end{enumerate}
\stfhSubIndex{I085}{Fuluoerbilian}{弗洛尔闭链（Floer cycle）}{Jibenbilian}{基本闭链（fundamental cycle）}
从 $f$ 的莫尔斯基本闭链经 $\mathcal L$ 传递得到的相应基本弗洛尔闭链具有如下良好
性质。该结果在 (KL03) 中对辛非球面流形证明，在 (Oh05b) 中推广到一般辛流形。

\begin{proposition}
\label{prop:22-canonical-fundamental-cycle}
设 $H$ 是通有的一周期哈密顿量，使得 $H_t$ 在 $x_-$ 处有唯一的非退化整体极小值，
而 $x_-$ 对所有 $t\in[0,1]$ 都是不动且欠扭的。设 $f:M\to\mathbb R$ 是以 $x_-$
为唯一整体极小点、并满足 $f(x_-)=0$ 的莫尔斯函数。则典范基本闭链具有表达式
\begin{equation}
 \alpha_H=[x_-,\widetilde x_-]+\beta\in CF(H),
 \label{eq:22-fundamental-cycle-expression}
\end{equation}
其中某条弗洛尔链 $\beta\in CF(H)$ 的层级满足
\begin{equation}
 \lambda_H(\beta)<\lambda_H([x_-,\widetilde x_-])
 =\int_0^1-H(t,x_-)\,dt.
 \label{eq:22-beta-level}
\end{equation}
特别地，$\alpha_H$ 的层级满足
\begin{align}
 \lambda_H(\alpha_H)
 &=\lambda_H([x_-,\widetilde x_-])
 \label{eq:22-fundamental-cycle-level}\\
 &=\int_0^1-H(t,x_-)\,dt
 =\int_0^1-\min H_t\,dt.\notag
\end{align}
\end{proposition}

\begin{proof}
证明基于下述引理（见 (KL03) 命题 5.2）。提醒读者，我们的符号约定与 (KL03) 不同。

\begin{lemma}
\label{lem:22-kerman-lalonde}
设 $H,f$ 如上。则对所有充分小的 $\epsilon>0$，定义为
\[
 G^H(t,x)=H(t,x_-)+\epsilon f(x)
\]
的函数 $G^H$ 满足
\[
 G^H(t,x_-)=H(t,x_-),\qquad G^H(t,x)\leq H(t,x)
\]
对所有 $(t,x)$ 成立，且等号只在 $x_-$ 处成立。
\end{lemma}

\begin{exercise}
\label{ex:22-prove-kerman-lalonde}
证明这个引理。
\end{exercise}

现在回到命题的证明。因为 $x_-$ 是 $H$ 与 $f$ 的欠扭不动极小点，所以有
康利--曾德尔指标
\[
 \mu_H([x_-,w_{x_-}])=\mu_{\epsilon f}([x_-,w_{x_-}])(=-n),
\]
从而模空间
$\mathcal M_{\mathcal L}([x_-,w_{x_-}],[x_-,w_{x_-}])$ 为零维。由于常值映射
$u(\tau,t)\equiv x_-$ 总在其中，故
\[
 \mathcal M_{\mathcal L}([x_-,w_{x_-}],[x_-,w_{x_-}])\neq\varnothing.
\]

任取
$u\in\mathcal M_{\mathcal L}([x_-,w_{x_-}],[x_-,w_{x_-}])$。注意，对线性同伦
\[
 \mathcal L:s\longmapsto(1-s)\epsilon f+sH
\]
的弗洛尔连续性方程，即使把同伦替换为
\[
 \mathcal L':s\longmapsto(1-s)G^H+sH,
\]
也保持不变。这是因为 $G^H$ 在 $\epsilon f$ 上增加的项 $H(t,x_-)$ 不依赖于
$x\in M$，因而
\[
 X_{\epsilon f}\equiv X_{G^H}.
\]
所以 $u$ 也是同伦 $\mathcal L'$ 下连续性方程的解。由此得到恒等式
\footnote{译者注：原书式 (22.2.16) 的第二个积分在括号内、外各印了一组
$dt\,d\tau$；译文按连续性能量恒等式只保留括号外的一组微分。}
\begin{align}
 \int\left\lVert\frac{\partial u}{\partial\tau}\right\rVert_{J^\rho}^2
 \,dt\,d\tau
 &=\mathcal A_{G^H}([x_-,w_{x_-}])
 -\mathcal A_H([x_-,w_{x_-}])\notag\\
 &\quad-\int_{-\infty}^{\infty}\rho'(\tau)
 \left(\int_0^1\bigl(H(t,u(\tau,t))-G^H(t,u(\tau,t))\bigr)\,dt\right)d\tau.
 \label{eq:22-continuation-energy-identity}
\end{align}
因为
\[
 \mathcal A_H([x_-,w_{x_-}])
 =\mathcal A_{G^H}([x_-,w_{x_-}])
 =\int_0^1-\min H\,dt
\]
且 $G^H\leq H$，故~\eqref{eq:22-continuation-energy-identity} 右端非正。因此
$\mathcal M_{\mathcal L}([x_-,w_{x_-}],[x_-,w_{x_-}])$ 只含常值解 $u\equiv x_-$。
特别地，这给出 $h_{\mathcal L}$ 的矩阵系数
\[
 \langle[x_-,w_{x_-}],h_{\mathcal L}([x_-,w_{x_-}])\rangle
 =\#\mathcal M_{\mathcal L}([x_-,w_{x_-}],[x_-,w_{x_-}])=1.
\]

现考虑 $\alpha_H$ 的任意其他生成元
\begin{equation}
 [z,w]\in\operatorname{supp}\alpha_H,
 \qquad [z,w]\neq[x_-,w_{x_-}].
 \label{eq:22-other-generator}
\end{equation}
由 $h_{\mathcal L}$ 与 $\alpha_H=h_{\mathcal L}(\alpha)$ 的定义，存在生成元
$[x,w_x]\in\operatorname{supp}\alpha$ 使得
\[
 \mathcal M_{\mathcal L}([x,w_x],[z,w])\neq\varnothing.
\]
于是对任意 $u\in\mathcal M_{\mathcal L}([x,w_x],[z,w])$，有恒等式
\[
 \begin{aligned}
 \mathcal A_H([z,w])-\mathcal A_{G^H}([x,w_x])
 &=-\int\left\lVert\frac{\partial u}{\partial\tau}\right\rVert_{J^\rho}^2
 \,dt\,d\tau\\
 &\quad-\int_{-\infty}^{\infty}\rho'(\tau)
 \left(\int_0^1(H(t,u(\tau,t))-G^H(t,u(\tau,t)))\,dt\right)d\tau.
 \end{aligned}
\]
由于 $-\lVert\partial u/\partial\tau\rVert_{J^\rho}^2\leq0$ 且 $G^H\leq H$，有
\begin{equation}
 \mathcal A_H([z,w])\leq\mathcal A_{G^H}([x,w_x]),
 \label{eq:22-action-comparison}
\end{equation}
等号只在 $u$ 为定常解时成立。分两种情况：$x=x_-$，以及
$[x_j,w_{x_j}]\in\operatorname{supp}\alpha$ 中 $x=x_j\neq x_-$。

第一种情形下，\emph{由假设} $[z,w]\neq[x_-,w_{x_-}]$，$u$ 不可能为常值，故
\eqref{eq:22-action-comparison} 中为严格不等式：
\begin{equation}
 \mathcal A_H([z,w])<\mathcal A_{G^H}([x_-,w_{x_-}]).
 \label{eq:22-strict-action-comparison}
\end{equation}
第二种情形下，某个 $x_j\neq x_-$ 满足~\eqref{eq:22-action-comparison}。注意，
\eqref{eq:22-morse-minimum-condition} 中的不等式 $f(x_-)<f(x_j)$ 等价于
\[
 \mathcal A_{G^H}([x_j,w_{x_j}])
 <\mathcal A_{G^H}([x_-,w_{x_-}]).
\]
它与~\eqref{eq:22-action-comparison} 一起再次给出
\eqref{eq:22-strict-action-comparison}。另一方面，因为
$G^H(t,x_-)=H(t,x_-)$，还有
\[
 \mathcal A_{G^H}([x_-,w_{x_-}])=\mathcal A_H([x_-,w_{x_-}]).
\]
综上，对任意
$[z,w]\in\operatorname{supp}\alpha_H$ 且 $[z,w]\neq[x_-,w_{x_-}]$，都有
\[
 \mathcal A_H([z,w])<\mathcal A_H([x_-,w_{x_-}])
 =\int_0^1-H(t,x_-)\,dt.
\]
命题得证。
\end{proof}

\subsection{自治哈密顿量的情形}

为证明极小性猜想，按照命题~\ref{prop:22-negative-length-minimizing} 给出的判据，只需证明
$\mathcal A_G([x_-,\widetilde x_-])=E^-(G)$ 与极小极大值 $\rho(G;1)$ 相等
（并类似证明 $\mathcal A_{\overline G}([x_+,\widetilde x_+])=E^-(\overline G)$
与 $\rho(\overline G;1)$ 相等）。后一事实是下述定理的直接推论。

\begin{theorem}
\label{thm:22-slow-fundamental-cycle-tight}
设 $G$ 是任意慢自治哈密顿量。则命题~\ref{prop:22-canonical-fundamental-cycle} 中构造的
典范基本闭链 $\alpha_G$ 是紧的，即
\[
 \rho(G;1)=\lambda_G(\alpha_G).
\]
特别地，$\rho(G;1)=E^-(G)$。类似地，
$\rho(\overline G;1)=E^-(\overline G)=E^+(G)$。
\end{theorem}

本节将证明这个定理。在开始证明前，先对其历史作一些说明。在以往多篇文献中，人们在
$G$ 满足下述条件或其变体时研究极小性猜想：
\begin{enumerate}
 \item $G$ 是慢的；
 \item 它有通有欠扭的极大点和极小点；
 \item 它的所有临界点在弗洛尔理论意义下非退化（即 $X_G$ 在每个临界点处的线性化流
 只有零点作为周期轨道）。
\end{enumerate}
例如参见 (MSl01), (En00), (Oh02), (KL03), (Oh05b)。条件 (2) 与 (3) 的出现，
其实是研究方法所造成的人为现象：该方法要使用某种形式的伪全纯曲线模空间，因而必须研究
模空间的横截性问题。通过这类研究可得的最强结果概括如下。

\begin{theorem}[定理 (Oh05b)]
\label{thm:22-oh-minimality}
设 $(M,\omega)$ 是任意闭辛流形，$G$ 是满足下列条件的自治哈密顿量：
\begin{enumerate}
 \item 它没有“周期为一”的非常值可缩周期轨道；
 \item 它有通有欠扭的极大点和极小点；
 \item 它的所有临界点在弗洛尔理论意义下非退化（即 $X_G$ 在每个临界点处的线性化流
 只有零点作为周期轨道）。
\end{enumerate}
则单参数子群 $\phi_G^t$ 对 $0\leq t\leq1$ 在其端点固定的同伦类中是长度极小的。
\end{theorem}

这个定理又是下述结果的直接推论。

\begin{theorem}
\label{thm:22-tight-under-twisted-cycle}
设 $G$ 是满足定理~\ref{thm:22-oh-minimality} 所列条件的任意慢自治哈密顿量。则命题
~\ref{prop:22-canonical-fundamental-cycle} 中构造的典范基本闭链 $\alpha_G$ 是紧的，即
\[
 \rho(G;1)=\lambda_G(\alpha_G)
 \bigl(=-G(x_-)=E^-(G)\bigr).
\]
\end{theorem}

\begin{proof}
注意，定理~\ref{thm:22-oh-minimality} 的条件特别蕴含 $G$ 非退化。固定一个不依赖时间的
$G$-正则殆复结构 $J_0$。

设 $\alpha$ 与典范基本弗洛尔闭链 $\alpha_G$ 同调，即
\begin{equation}
 \alpha=\alpha_G+\partial_G(\gamma)
 \label{eq:22-homologous-fundamental-cycle}
\end{equation}
对某条弗洛尔链 $\gamma\in CF_*(G)$ 成立。当 $G$ 自治且 $J\equiv J_0$ 不依赖于
$t$ 时，不存在落在 $[x_-,\widetilde x_-]$ 的非定常且不依赖于 $t$ 的
$\mathcal A_G$ 流线；因为任何这样的流线都来自 $G$ 的负莫尔斯梯度流，而 $x_-$
是 $G$ 的极小点。因此，任何落在 $[x_-,\widetilde x_-]$ 的非定常弗洛尔流线 $u$
必依赖于 $t$。由于假设 $G$ 不存在周期为一的非常值可缩周期轨道，
$\mathcal A_G$ 的任一临界点都形如
\[
 [x,w],\qquad x\in\operatorname{Crit}G.
\]

设 $u$ 是从 $[x,w]$（$x\in\operatorname{Crit}G$）出发、满足
\begin{equation}
 \mu([x,w])-\mu([x_-,\widetilde x_-])=1
 \label{eq:22-index-one-trajectory}
\end{equation}
的流线，并以
$\widehat{\mathcal M}_{(G,J_0)}([x,w],[x_-,\widetilde x_-])$ 表示相应的连接流线
弗洛尔模空间。一般指标公式 (18.4.24) 给出
\begin{equation}
 \mu([x,w])=\mu([x,w_x])-2c_1([w]).
 \label{eq:22-index-capping-formula}
\end{equation}
分别考虑 $c_1([w])=0$ 与 $c_1([w])\neq0$ 两种情形。若
$c_1([w])\neq0$，由~\eqref{eq:22-index-one-trajectory} 与
~\eqref{eq:22-index-capping-formula} 得 $x\neq x_-$。这意味着任意这样的流线都带有
（局部）自由的 $S^1$-作用（见第 19.8 节），即模空间
\[
 \mathcal M_{(G,J_0)}([x,w],[x_-,\widetilde x_-])
 =\widehat{\mathcal M}_{(G,J_0)}([x,w],[x_-,\widetilde x_-])/\mathbb R
\]
及其稳定映射紧化
$\overline{\mathcal M}_{(G,J_0)}([x,w],[x_-,\widetilde x_-])$ 都带有\emph{无不动点}的局部
自由 $S^1$-作用。于是由 (FHS95) 的 $S^1$-等变横截性定理，对适当选择的自治
$J_0$，模空间
$\mathcal M_{(G,J_0)}([x,w],[x_-,\widetilde x_-])$ 为空。这是因为在假设
~\eqref{eq:22-index-one-trajectory} 下，其商空间的虚维数为 $-1$。关于这一
$S^1$-不变正则化过程的更多说明，参见 (FHS95)。

现在考虑 $c_1([w])=0$ 的情形。首先，
~\eqref{eq:22-index-one-trajectory} 与~\eqref{eq:22-index-capping-formula}
蕴含 $x\neq x_-$。另一方面，若 $x\neq x_-$，上述同一论证表明扰动后的模空间为空。

因此，不存在任何落在 $[x_-,\widetilde x_-]$ 的指标 $1$ 流线。所以不论如何选择
$\gamma$，~\eqref{eq:22-homologous-fundamental-cycle} 中的 $\partial_G(\gamma)$
都不能消去闭链
\[
 \alpha_G=[x_-,\widetilde x_-]+\partial_G(\gamma)
\]
中的项 $[x_-,\widetilde x_-]$，故由层级 $\lambda_G$ 的定义有
\[
 \lambda_G(\alpha)\geq\lambda_G([x_-,\widetilde x_-]).
\]
结合~\eqref{eq:22-fundamental-cycle-level} 即完成证明。
\end{proof}

为把定理~\ref{thm:22-oh-minimality} 推广到极小性猜想中的一般慢哈密顿量（不必欠扭），
我们把下述逼近结果与 $\rho(\mathord\cdot;1)$ 的 $C^0$-连续性结合使用。厄舍在
(Ush10a) 定理 4.5 中证明了这一结果；它改进了 (Sch06) 定理 1.3，并在逼近过程中
增补了非退化性。

\begin{theorem}
\label{thm:22-slow-approximation-preview}
设 $H:M\to\mathbb R$ 是任意慢哈密顿量。则对任意 $\epsilon>0$，存在光滑函数
$K:M\to\mathbb R$，使 $\lVert K-H\rVert_{C^0}<\epsilon$，且 $K$ 是慢的、欠扭的
莫尔斯函数。
\end{theorem}

\begin{proof}
证明见下一小节，即第 22.2.5 小节。
\end{proof}

最后总结极小性猜想的证明。

\begin{theorem}[厄舍 (Ush10a)]
\label{thm:22-usher-minimality}
设 $(M,\omega)$ 是任意闭辛流形，$G$ 是任意慢自治哈密顿量。则单参数群
$\phi_G^t$ 对 $0\leq t\leq1$ 在其端点固定的同伦类中是长度极小的。
\end{theorem}

\begin{proof}
这是命题~\ref{prop:22-negative-length-minimizing} 与定理
~\ref{thm:22-slow-fundamental-cycle-tight} 的直接推论。更精确地，定理
~\ref{thm:22-slow-fundamental-cycle-tight} 蕴含 $\rho(G;1)=E^-(G)$，再由命题
~\ref{prop:22-negative-length-minimizing} 得 $G$ 使负霍费尔长度极小。因为 $G$
慢时 $\overline G$ 也慢，还有
$\rho(\overline G;1)=E^-(\overline G)=E^+(G)$。因此 $\overline G$ 也使负霍费尔
长度极小，这等价于 $G$ 使正长度极小。两者结合显然说明 $G$ 使霍费尔长度极小。
\end{proof}

\subsection{慢哈密顿量的逼近定理}

本小节旨在逐字依照 (Ush10a) 中厄舍的证明，给出定理
~\ref{thm:22-slow-approximation-preview} 的证明。

\stfhIndex{I200}{Yueke}{约克（Yorke）}
先从约克 (Yo69) 得到的一般常微分方程周期估计开始。设
$F:\Omega\to\mathbb R^n$ 是区域 $\Omega\subset\mathbb R^n$ 上的利普希茨
向量场，以 $\lVert\mathord\cdot\rVert$ 表示欧氏范数。约克 (Yo69)
证明了如下结果。（证明细节请读者参阅原始论文。）

\begin{theorem}
\label{thm:22-yorke-period-estimate}
设 $F$ 是利普希茨常数为 $L>0$ 的利普希茨映射，即
\begin{equation}
 \lVert F(x_1)-F(x_2)\rVert\leq L\lVert x_1-x_2\rVert
 \qquad\text{对每个 }x_1,x_2\in\Omega.
 \label{eq:22-lipschitz-vector-field}
\end{equation}
若 $x=x(t)$ 是 $\dot x=F(x)$ 的周期为 $p$ 的非常值周期解，则
$p\geq2\pi/L$。
\end{theorem}

涩谷（未发表工作）更早证明过一条类似定理，所得较弱结论为 $p\geq2/L$；对这里的
目的而言，这其实已经足够。

这里再次回忆定理~\ref{thm:22-slow-approximation-preview} 的陈述。

\begin{theorem}
\label{thm:22-slow-approximation}
设 $H:M\to\mathbb R$ 是任意慢哈密顿量。则对任意 $\epsilon>0$，存在光滑函数
$K:M\to\mathbb R$，使 $\lVert K-H\rVert_{C^0}<\epsilon$，且 $K$ 是慢的、通有
欠扭的莫尔斯函数。
\footnote{译者注：原书定理 22.2.17 的预告版本称逼近函数为“欠扭”，本定理重述时则称
“通有欠扭”；译文分别照录。}
\end{theorem}

\begin{proof}
选择由与 $\omega$ 相容的殆复结构诱导的 $M$ 上背景黎曼度量，使得
$X_H=J\nabla H$，特别地，对所有 $p$ 都有 $|X_H(p)|=|\nabla H(p)|$。令
\[
 F_n=\{p\in M\mid |\nabla H(p)|\leq1/n\}.
\]
注意，若 $s\in\bigcap_{n=1}^\infty H(F_n)$，则存在 $p_n\in M$ 使
$H(p_n)=s$ 且 $|\nabla H(p_n)|\leq1/n$；由 $M$ 的紧性以及 $H,\nabla H$ 的连续性，
存在 $p\in M$ 满足 $H(p)=s$ 且 $\nabla H(p)=0$。故
$\bigcap_{n=1}^\infty H(F_n)$ 恰为 $H$ 的临界值集合，由萨德定理它的勒贝格
测度为零。因为 $H$ 的像有有限勒贝格测度，故当 $n\to\infty$ 时
$m(H(F_n))\to0$。

给定 $\zeta>0$。存在 $N$ 使 $m(H(F_N))<\zeta/2$。选择光滑函数
$\psi:M\to\mathbb R$，使
\[
 0\leq\psi\leq1,\qquad \psi|_{F_{2N}}=1,
 \qquad\operatorname{supp}(\psi)\subset F_N.
\]
因为 $F_N$ 从而 $H(F_N)\subset\mathbb R$ 紧，所以存在闭包两两不交的开区间
$I_1,\ldots,I_k$，使
\[
 H(F_N)\subset\bigcup_{j=1}^k I_j,
 \qquad\sum_{j=1}^k m(I_j)<\frac{3\zeta}{4}.
\]
再选闭包两两不交的开区间 $I'_1,\ldots,I'_k$，使对每个 $j$ 都有
$I_j\Subset I'_j$ 且 $\sum_{j=1}^k m(I'_j)<\zeta$。记
\[
 S=\bigcup_{j=1}^k I_j,
 \qquad S'=\bigcup_{j=1}^k I'_j.
\]
令
\[
 B=\sup_{p\in M}|\nabla(\nabla H)(p)|,
\]
并取一个稍后还要进一步指定的小数 $\eta>0$；特别要求 $\eta<1$。选择
$C^\infty$ 函数
$f:[\min H,\max H]\to[\min H,\max H]$ 满足：
\begin{enumerate}
 \item $f(\min H)=\min H$；
 \item 对所有 $s$，$0<f'(s)\leq1$；
 \item 若 $s\notin S'$，则 $f'(s)=1$；
 \item 若 $s\in S$，则 $f'(s)=\eta/2$。
\end{enumerate}
因为 $m(S')<\zeta$，对所有 $s\in[\min H,\max H]$ 都有
$0\leq s-f(s)<\zeta$。特别地，$\lVert f\circ H-H\rVert_{C^0}<\zeta$，且此估计
与 $\eta$ 无关。

若 $q\in F_N$（由 $f$ 的构造，这特别蕴含 $f''(H(q))=0$），则
\[
 |\nabla(f\circ H)(q)|=f'(H(q))|\nabla H(q)|\leq\frac{\eta}{2N}
\]
且
\[
 |\nabla(\nabla(f\circ H))(q)|
 =\frac{\eta}{2}|\nabla(\nabla H)(q)|\leq\frac{B\eta}{2}.
\]
故
\[
 \lVert X_{f\circ H}\rVert_{C^1}
 \leq\frac{\eta}{2}+\frac{B\eta}{2}.
\]
回忆上述函数 $\psi:M\to\mathbb R$（其支集包含于 $F_N$），于是
\footnote{译者注：此处 $\psi X_{f\circ H}$ 表示函数 $\psi$ 与向量场
$X_{f\circ H}$ 的逐点乘积；源书即印作并列乘积，无复合圈。}
\[
 \lVert\psi X_{f\circ H}\rVert_{C^1}
 \leq\left(\frac{\eta}{2}+\frac{B\eta}{2}\right)
 \lVert\psi\rVert_{C^1}.
\]
由定理~\ref{thm:22-yorke-period-estimate} 的估计（以及惠特尼嵌入定理），存在
$\eta_0>0$，使 $M$ 上任意满足
$\lVert V\rVert_{C^1}\leq\eta_0(B+1)\lVert\psi\rVert_{C^1}$ 的向量场 $V$
都没有周期小于或等于一的非常值周期轨道。

于是取 $0<\eta<\eta_0$，再取 $M$ 上一列莫尔斯函数
$\{K_m\}_{m=1}^\infty$，使 $K_m$ 在 $C^2$-范数下趋于 $f\circ H$。则
$X_{K_m}\to X_{f\circ H}$ 于 $C^1$-范数，因而
$\psi X_{K_m}\to\psi X_{f\circ H}$ 于 $C^1$-范数。所以当 $m$ 足够大，使
\[
 \lVert\psi X_{f\circ H}-\psi X_{K_m}\rVert_{C^1}
 <\frac{(B+1)\eta}{2}\lVert\psi\rVert_{C^1},
\]
向量场 $\psi X_{K_m}$ 没有周期小于或等于一的非常值周期轨道。

现在，因为 $H$ 慢且 $0<f'\leq1$，容易看出 $f\circ H$ 也慢。因此
$X_{f\circ H}$ 的所有周期小于或等于一的可缩周期轨道，都是 $f\circ H$ 临界点处
的常值轨道；特别地，它们全都包含在 $F_{2N}$ 的内部。阿尔泽拉--阿斯科利定理的一个
简单应用于是表明：若 $m$ 充分大，$X_{K_m}$ 的任意周期小于或等于一的可缩周期轨道
必包含在 $F_{2N}$ 中。因此（由于 $\psi|_{F_{2N}}=1$），对充分大的 $m$，
$X_{K_m}$ 的周期小于或等于一的可缩周期轨道其实也是 $\psi X_{K_m}$ 的周期轨道，
从而由上文可知只能是常值轨道。所以充分大的 $m$ 使 $K_m$ 为慢的。

此外，$K_m\to f\circ H$ 于 $C^2$ 拓扑蕴含：当 $m$ 充分大时，$K_m$ 的所有临界点
$p$ 都包含于 $F_N$；在任一这样的 $p$ 处，
\[
 |\nabla(\nabla(f\circ H))(p)|<B\eta/2,
\]
所以当 $m$ 充分大时 $|\nabla(\nabla K_m)(p)|<B\eta$。因此，只要起初把 $\eta$
取成小于 $2\pi/B$，$K_m$ 在任一此类临界点处的线性化流 $\phi_{K_m}^t$ 就没有周期
小于或等于 $1$ 的非常值周期轨道。于是对充分大的 $m_0$，$K_{m_0}$ 是慢的、平坦的
莫尔斯函数，且
$\lVert K_{m_0}-f\circ H\rVert_{C^2}<\eta$。令 $\zeta=\epsilon/2$，并取
$K=K_{m_0}$，再结合 $\lVert f\circ H-H\rVert_{C^0}<\eta$，即证定理。
\footnote{译者注：原证明取 $\zeta=\epsilon/2$，但对
$\lVert K_{m_0}-f\circ H\rVert_{C^2}<\eta$ 及
$\lVert f\circ H-H\rVert_{C^0}<\eta$ 未再明确要求 $\eta<\epsilon/2$；此外，此处结论
称 $K_{m_0}$ “平坦”，而定理陈述称“通有欠扭”。
译文保留论证结构并提示这两处缺口。}
\end{proof}
